\section{Conclusion}
\label{sec:conclusion}

A mandatory credit-risk disclosure moved uninsured funding at community
banks.  The CECL Day-One adjustment---an accounting charge whose
magnitude was fixed by past underwriting---produced a within-bank shift
of time deposits from uninsured to insured balances in proportion to
the disclosed number, with flat five-year pre-adoption trajectories.
The timing identifies who acted on what information: banks, which knew
their charges privately from the 2022 parallel runs, pre-funded with
insured brokered deposits and elected capital relief before adoption;
depositors moved only when the number became public.  Discipline
operated as quantity exit at the rollover margin---high-CECL banks
never repriced retail deposits, in the stress window or after---and
banks repurchased the lost coverage in wholesale insured markets, at
the cost of higher interest expense and lower profitability.

The policy implication is narrow and conditional.  Legislation that
would expand deposit insurance coverage for business accounts, such as
the Main Street Depositor Protection Act, would convert uninsured
funding of the kind I study into insured balances that carry no
incentive to monitor.  My estimates identify the marginal cost of that
substitution on one channel only---the funding-cost and profitability
consequences of a mandatory disclosure at community banks in
2023---and do not measure the run-prevention benefits of broader
coverage or the general-equilibrium response of bank risk-taking.
What the evidence does support is a weaker statement: proposals to
expand coverage should be assessed for the monitoring and pricing
discipline they displace as well as for the stability gains they
deliver.
